\chapter{Causal Residue}

\section{What Appearance Doesn't Tell You}

Consider two courtyards, identical in every respect that appearance can capture. Both have a stone bench near the far wall, weathered in the same way, positioned identically, indistinguishable to any observation reading z alone. In the first courtyard, the bench has stood in that spot since the wall itself was built. In the second, the bench was moved there yesterday, replacing an older bench that stood a few feet to the left. Appearance, as this book has developed it, cannot tell these two courtyards apart. \(z_t(x)\) reports what a point currently looks like, and at this moment, both points look exactly the same.

The two courtyards are not, however, equivalent, and something in this framework needs to be able to say so. If a visitor arrives asking where the old bench went, the first courtyard has no answer to give because there was no old bench; the second courtyard owes an account, because something displaced. If a structural inspector wants to know whether the ground beneath the bench has borne this particular load for decades or for a single day, appearance alone cannot answer, but the difference matters for what is admissible to happen next. Two states of the field can be visually and even structurally identical in the present while permitting entirely different futures, because what is admissible going forward depends not only on what a point currently is but on what sequence of events produced it. This is what appearance, on its own, cannot supply, and it is the reason the five-component decomposition set aside a slot for it as early as Chapter 7 without yet explaining why.

\section{Residue as Compressed History}

The obvious response is to propose that the field simply keep a complete log: every event, at every point, timestamped and retained indefinitely, so that any question about history can always be answered by consulting the record directly. This response should be resisted, and resisting it is the substantive content of this section.

A complete, unbounded log is not what \(R_t\) denotes in this book, for two reasons that turn out to be related. The practical reason is the obvious one: a field that retained every event at every point without ever discarding anything would grow without bound, and a persistent world that cannot be maintained indefinitely without unbounded storage has not solved the persistence problem, it has merely deferred its cost. The deeper reason is that most of what a full log would contain is not actually needed to determine what remains admissible. What matters about the courtyard's history is not the complete sequence of everything that ever happened there, but whatever subset of that sequence continues to constrain the present: that a bench was moved, roughly when, and what displaced it, not the precise angle of every hand that touched it in transit.

Residue, then, is compressed history: a summary trace, retained at each point of the field, sufficient to determine what remains admissible without retaining everything that produced it. This compression is not merely an engineering convenience bolted on afterward. It is itself an event, in the same sense that observation was shown in Chapters 8 and 9 to be an event rather than a passive read. When two histories are merged, or when an old event's specific detail is folded into a coarser summary because nothing downstream still depends on that detail, information is genuinely and irreversibly lost, in exactly the sense that merging two branches of a system's history always discards whatever was incompatible between them. Residue is what remains after this discarding, and it is worth being honest that discarding has occurred rather than treating \(R_t\) as though it were a lossless archive that merely happens to be summarized for convenience.

\section{How Residue Accumulates}

Residue is not fixed at the moment a point is first observed, the way this book's earlier chapters might suggest by analogy with appearance. It accumulates, receiving a contribution each time an event actually occurs at or near that point, whether or not the point is currently being observed. Schematically,

\[ R_{t+1}(x) = R_t(x) \oplus \mathrm{event}_t(x), \]

where ⊕ denotes whatever accumulation operator combines existing residue with a new event's trace. This book will not attempt, at this stage, to pin down ⊕ precisely; unlike z's update rule, which is a genuine weighted average admitting a clean closed form, residue accumulation is closer to a structural merge, more like combining two branches of a causal graph than blending two numbers, and its full algebraic treatment belongs to Part III and Part IV, once constraint projection and the causal-DAG structure of multi-agent updates have been properly developed. What matters here is only the shape of the claim: residue grows by accumulation rather than being fixed once and held constant, and each accumulation step is itself subject to the same discipline against unearned change that governed entropy and appearance in the two preceding chapters. A point's residue does not drift on its own; it only accumulates in response to something that actually happened.

\section{Residue and Affordance}

Chapter 10 introduced a learning rule for affordances, in which the weight connecting a cue to an affordance strengthens with repeated, consistent encounters and weakens or fails to consolidate when encounters disagree. It is worth noticing now, with residue properly in view, that this learning rule was already describing a form of residue accumulation without naming it as such. The weight \(w_{ij}\) at a given location is itself a compressed trace of that location's history of activating affordance j in response to cue i: not a full log of every encounter, but a summary sufficient to determine how strongly that association should hold going forward. Residue and learned affordance are not two unrelated mechanisms that happen to share a family resemblance. Affordance learning is a special case of residue accumulation, restricted to the particular kind of event this book calls an encounter, and this closes a loop between Chapter 10 and this chapter that was left open when affordances were first introduced as though they arose fresh from cues with no reference to what had come before.

\section{Residue as a Constraint on Future Invention}

This chapter closes by returning to the invention-to-memory principle from Chapter 9 and strengthening it. That principle stated that an already-settled appearance cannot be casually overwritten by a fresh candidate. It is now possible to say what an unsettled point's candidate must actually be checked against, and the answer is not only the point's own prior appearance, which for a genuinely unresolved point may not exist yet, but the point's accumulated residue. A courtyard whose residue records that a bench was recently moved constrains what invention is permitted to produce there: an observer generating a first appearance for the newly bare patch of ground the bench used to occupy is not free to invent anything at all, because residue has already ruled out large classes of otherwise-plausible content, ground that has clearly been recently disturbed does not admit an invented appearance of undisturbed, decades-settled paving. Residue, in other words, does not merely record what happened for the sake of answering questions later. It actively narrows the space of admissible invention going forward, and in this sense it is doing for entropy exactly what appearance was shown to do for observation in Chapter 9: giving Chapter 8's entropy constraint teeth, by specifying not only that unearned certainty is forbidden but which specific certainties remain earnable given what has already occurred.

\section{Closing Part II}

This chapter completes the account this book set out to give in Chapter 6. The world is a persistent field, \(M_t\) : X → \(\mathbb{R}^d\), and that field's value at every point decomposes into five components, each answering a question the others cannot answer for it: Φ, how stable a region is; v, where it is moving; S, how much about it remains genuinely open; R, what has happened there; and z, what it currently looks like. Affordance, the space of actions a given capability can take at a given point, is not a sixth component stored alongside these five, but a relation arising at their intersection with an agent's own capability, triggered by cues, learned through repetition, and coherent or incoherent across an agent's encounters in a sense this book has now made precise.

Nothing in the last six chapters has said how any of this changes. Φ, v, S, R, and z have each been described in terms of what they hold and what constrains their update, but the actual mechanism by which a proposed change becomes a committed one, the machinery hinted at as early as Chapter 6's five-stage cycle and touched again in the invention-to-memory principle, has been assumed rather than built. Part III takes up that construction directly. It begins where this book's core cycle diagram always pointed: with \(\Delta\psi\), the proposal, and the question of where a proposal like that actually comes from.
